\chapter{Camera acquisition}\label{module:web-camera}
\Q{Evidence}{Current implementation boundary}
\path{Software/DemoCode/web/camera.py}. Exact file hashes are in the accompanying source manifest. Names, declarations and source locations below come from the captured source, not a runtime import.
\Q{Purpose}{Input-to-output reading flow}
Acquire a still image or preview frame under a shared camera lock, then pass captured pixels to the page-processing stage. Preview lifecycle and still capture share the device.


\Q{Evidence}{Selected current module settings}
These are declarations or expressions, not recovered environment values. Source comments do not establish a measured optimum. Larger glyph coordinate tables are designed data in the bundled source.
\begin{lstlisting}[style=plain]
SCRIPT_DIR = Path(__file__).resolve().parent
CAPTURE_DIR = paths.Cache / 'CameraCaptures'
CAMERA_INDEX = int(os.environ.get('CAMERA_INDEX', '1'))
WARMUP_FRAMES = 5
\end{lstlisting}
\section*{Stage: Camera capture and image orientation}
\section{Method background: Camera capture and image orientation}\label{method:camera}
\Q{Purpose}{The role in this system}
Acquire a frame and put its pixel orientation in agreement with the browser and segmenter.


\Q{Justification}{The reason for this construction}
An existing preview handle avoids competing camera opens; warm-up frames allow acquisition to settle. EXIF transpose applies metadata to pixels rather than assuming the viewer appearance equals array orientation.


\Q{Alternatives}{What changes if another method is used}
Direct V4L2 access offers lower-level controls; a CSI camera uses a different acquisition stack. Reading one immediate frame may capture an unsettled exposure/focus state.


\Q{Library}{Implementation speed and accuracy}
OpenCV VideoCapture and Pillow ImageOps.exif\_transpose are already used. Their presence concerns acquisition/decoding and orientation, not a claim that OpenCV implements the custom segmentation. Sensor claims need actual negotiated settings. \cite{opencvcamera,pillowops}.


\section{\_GetLock}\label{fn:web-camera:-GetLock}
\begin{lstlisting}[style=sig]
def _GetLock()
\end{lstlisting}
\Q{Evidence}{Source and method link}
\path{Software/DemoCode/web/camera.py} line 39. Implementation evidence: \cite{code2026}. Related method, alternatives and library discussion: Section~\ref{method:camera}.


\Q{Purpose}{What this contributes}
Acquire a frame and put its pixel orientation in agreement with the browser and segmenter.


\Q{Method}{Current implementation}
Acquire a frame and put its pixel orientation in agreement with the browser and segmenter. The excerpt below is from the current definition; the complete function is in the source snapshot. Its constants and branches, rather than an older signature, govern this path.


\begin{lstlisting}[style=plain]
def _GetLock():
    global _previewLock
    if _previewLock is None:
        import threading
        _previewLock = threading.Lock()
    return _previewLock
\end{lstlisting}
\Q{Justification}{Interpret the method rather than infer a benchmark}
An existing preview handle avoids competing camera opens; warm-up frames allow acquisition to settle. EXIF transpose applies metadata to pixels rather than assuming the viewer appearance equals array orientation. \cite{opencvcamera,pillowops}.


\section{PreviewStart}\label{fn:web-camera:PreviewStart}
\begin{lstlisting}[style=sig]
def PreviewStart(index=None)
\end{lstlisting}
\Q{Evidence}{Source and method link}
\path{Software/DemoCode/web/camera.py} line 52. Implementation evidence: \cite{code2026}. Related method, alternatives and library discussion: Section~\ref{method:camera}.


\Q{Purpose}{What this contributes}
Opens (and keeps open) the camera for live preview streaming. Safe to call repeatedly -- a no-op if already open. Kept separate from CaptureImage() so a short-lived capture doesn't leave the sensor powered on between uses -- the camera only runs while someone's actually watching the live feed to frame the shot and let autofocus settle.


\Q{Method}{Current implementation}
Acquire a frame and put its pixel orientation in agreement with the browser and segmenter. The excerpt below is from the current definition; the complete function is in the source snapshot. Its constants and branches, rather than an older signature, govern this path.


\begin{lstlisting}[style=plain]
def PreviewStart(index=None):
    import cv2
    global _previewCap
    with _GetLock():
        if _previewCap is not None:
            return
        idx = CAMERA_INDEX if index is None else index
        cap = cv2.VideoCapture(idx)
        if not cap.isOpened():
            cap.release()
            raise RuntimeError(f"Could not open camera at index {idx} for preview.")
        _previewCap = cap
\end{lstlisting}
\Q{Justification}{Interpret the method rather than infer a benchmark}
An existing preview handle avoids competing camera opens; warm-up frames allow acquisition to settle. EXIF transpose applies metadata to pixels rather than assuming the viewer appearance equals array orientation. \cite{opencvcamera,pillowops}.


\section{PreviewStop}\label{fn:web-camera:PreviewStop}
\begin{lstlisting}[style=sig]
def PreviewStop()
\end{lstlisting}
\Q{Evidence}{Source and method link}
\path{Software/DemoCode/web/camera.py} line 68. Implementation evidence: \cite{code2026}. Related method, alternatives and library discussion: Section~\ref{method:camera}.


\Q{Purpose}{What this contributes}
Releases the preview camera handle. Always call this once you're done framing the shot -- see the module comment's note on heat.


\Q{Method}{Current implementation}
Acquire a frame and put its pixel orientation in agreement with the browser and segmenter. The excerpt below is from the current definition; the complete function is in the source snapshot. Its constants and branches, rather than an older signature, govern this path.


\begin{lstlisting}[style=plain]
def PreviewStop():
    global _previewCap
    with _GetLock():
        if _previewCap is not None:
            _previewCap.release()
            _previewCap = None
\end{lstlisting}
\Q{Justification}{Interpret the method rather than infer a benchmark}
An existing preview handle avoids competing camera opens; warm-up frames allow acquisition to settle. EXIF transpose applies metadata to pixels rather than assuming the viewer appearance equals array orientation. \cite{opencvcamera,pillowops}.


\section{PreviewActive}\label{fn:web-camera:PreviewActive}
\begin{lstlisting}[style=sig]
def PreviewActive()
\end{lstlisting}
\Q{Evidence}{Source and method link}
\path{Software/DemoCode/web/camera.py} line 76. Implementation evidence: \cite{code2026}. Related method, alternatives and library discussion: Section~\ref{method:camera}.


\Q{Purpose}{What this contributes}
Acquire a frame and put its pixel orientation in agreement with the browser and segmenter.


\Q{Method}{Current implementation}
Acquire a frame and put its pixel orientation in agreement with the browser and segmenter. The excerpt below is from the current definition; the complete function is in the source snapshot. Its constants and branches, rather than an older signature, govern this path.


\begin{lstlisting}[style=plain]
def PreviewActive():
    return _previewCap is not None
\end{lstlisting}
\Q{Justification}{Interpret the method rather than infer a benchmark}
An existing preview handle avoids competing camera opens; warm-up frames allow acquisition to settle. EXIF transpose applies metadata to pixels rather than assuming the viewer appearance equals array orientation. \cite{opencvcamera,pillowops}.


\section{PreviewFrameJpeg}\label{fn:web-camera:PreviewFrameJpeg}
\begin{lstlisting}[style=sig]
def PreviewFrameJpeg()
\end{lstlisting}
\Q{Evidence}{Source and method link}
\path{Software/DemoCode/web/camera.py} line 83. Implementation evidence: \cite{code2026}. Related method, alternatives and library discussion: Section~\ref{method:camera}.


\Q{Purpose}{What this contributes}
Grabs one frame from the already-open preview camera and returns JPEG-encoded bytes, for an MJPEG-style streaming endpoint. Raises RuntimeError if PreviewStart() hasn't been called.


\Q{Method}{Current implementation}
Acquire a frame and put its pixel orientation in agreement with the browser and segmenter. The excerpt below is from the current definition; the complete function is in the source snapshot. Its constants and branches, rather than an older signature, govern this path.


\begin{lstlisting}[style=plain]
def PreviewFrameJpeg():
    import cv2
    with _GetLock():
        if _previewCap is None:
            raise RuntimeError("Preview not started -- call preview_start() first.")
        ok,frame = _previewCap.read()
    if not ok:
        raise RuntimeError("Failed to read a preview frame.")
    ok,buf = cv2.imencode(".jpg",frame)
    if not ok:
        raise RuntimeError("Failed to JPEG-encode preview frame.")
    return buf.tobytes()
\end{lstlisting}
\Q{Justification}{Interpret the method rather than infer a benchmark}
An existing preview handle avoids competing camera opens; warm-up frames allow acquisition to settle. EXIF transpose applies metadata to pixels rather than assuming the viewer appearance equals array orientation. \cite{opencvcamera,pillowops}.


\section{CaptureImage}\label{fn:web-camera:CaptureImage}
\begin{lstlisting}[style=sig]
def CaptureImage(timeout=10, index=None)
\end{lstlisting}
\Q{Evidence}{Source and method link}
\path{Software/DemoCode/web/camera.py} line 105. Implementation evidence: \cite{code2026}. Related method, alternatives and library discussion: Section~\ref{method:camera}.


\Q{Purpose}{What this contributes}
Captures one still image, returns its Path. Raises RuntimeError if no camera could be opened/read -- callers (Server.py) should catch this and report it as a clean error, not a crash. if a live preview is currently open (PreviewStart()), this grabs the frame from THAT handle instead of opening a second one -- most webcams only allow one process/handle to hold the device at a time, and this is also the intended flow: watch the live preview until autofocus settles on the page, then capture "the good image" from what's already showing.


\Q{Method}{Current implementation}
Acquire a frame and put its pixel orientation in agreement with the browser and segmenter. The excerpt below is from the current definition; the complete function is in the source snapshot. Its constants and branches, rather than an older signature, govern this path.


\begin{lstlisting}[style=plain]
def CaptureImage(timeout=10,index=None):
    import cv2

    with _GetLock():
        if _previewCap is not None:
            ok,frame = _previewCap.read()
            if not ok:
                raise RuntimeError("Failed to read a frame from the open preview.")
            outPath = CAPTURE_DIR/f"capture_{int(time.time()*1000)}.jpg"
            cv2.imwrite(str(outPath),frame)
            return outPath

    idx = CAMERA_INDEX if index is None else index
    cap = cv2.VideoCapture(idx)
    if not cap.isOpened():
        cap.release()
        raise RuntimeError(
            f"Could not open camera at index {idx}. Is it plugged in, and "
            "not already in use by another program (close any other app "
            "using the webcam, e.g. a video call)?"
        )

    deadline = time.time()+timeout
    frame = None
    try:
        for _ in range(WARMUP_FRAMES):
            if time.time()>deadline:
                break
            ok,frame = cap.read()
            if not ok:
                frame = None
    finally:
        cap.release()

    if frame is None:
        raise RuntimeError(
            f"Camera at index {idx} opened but produced no readable frame "
            "within the timeout."
        )

    outPath = CAPTURE_DIR/f"capture_{int(time.time()*1000)}.jpg"
    cv2.imwrite(str(outPath),frame)
    if not outPath.exists():
        raise RuntimeError(f"cv2.imwrite reported no error but wrote no file: {outPath}")
    return outPath
\end{lstlisting}
\Q{Justification}{Interpret the method rather than infer a benchmark}
An existing preview handle avoids competing camera opens; warm-up frames allow acquisition to settle. EXIF transpose applies metadata to pixels rather than assuming the viewer appearance equals array orientation. \cite{opencvcamera,pillowops}.

